\documentclass{article}
\usepackage{amsmath,amssymb,amsthm}
\usepackage{algorithm,algorithmic}
\usepackage{hyperref}
\usepackage{enumitem}
\newtheorem{theorem}{Theorem}
\newtheorem{lemma}{Lemma}
\newtheorem{assumption}{Assumption}
\newcommand{\prox}{\operatorname{prox}}
\newcommand{\R}{\mathbb{R}}
\newcommand{\E}{\mathbb{E}}
\newcommand{\norm}[1]{\left\lVert#1\right\rVert}
\begin{document}

\section*{Algorithm Name}
\textbf{Accelerated Proximal Gradient with Momentum (APG-M)}  

We consider the composite convex problem  

\[
\min_{x\in\R^{d}} \; F(x)\;:=\;f(x)+g(x),
\]

where  

\begin{itemize}
  \item $f:\R^{d}\to\R$ is convex and differentiable with $L$‑Lipschitz continuous gradient,
  \item $g:\R^{d}\to\R\cup\{+\infty\}$ is convex, possibly non‑smooth, and its proximal operator
        \[
          \prox_{\eta g}(v):=\arg\min_{u\in\R^{d}}\Big\{ g(u)+\frac{1}{2\eta}\norm{u-v}^{2}\Big\}
        \]
        can be evaluated efficiently.
\end{itemize}

The algorithm augments the classical proximal gradient step with a Nesterov‑type momentum on the \emph{gradient input}.

\section*{Mathematical Formulation}
Given a stepsize $\eta\in(0,1/L]$ and a sequence of momentum parameters $\{\beta_k\}_{k\ge 0}\subset[0,1)$, the iterates are generated as follows:

\[
\boxed{
\begin{aligned}
y_k &:= x_k + \beta_k (x_k-x_{k-1}), \qquad (k\ge 0,\;x_{-1}=x_0)\\[4pt]
x_{k+1} &:= \prox_{\eta g}\!\bigl( y_k - \eta \nabla f(y_k) \bigr).
\end{aligned}}
\tag{APG-M}
\]

Typical choices (FISTA) are  

\[
t_{k+1}= \frac{1+\sqrt{1+4t_k^{2}}}{2},\qquad 
\beta_k = \frac{t_k-1}{t_{k+1}},\qquad t_0=1 .
\]

\section*{Pseudocode}
\begin{algorithm}[H]
\caption{Accelerated Proximal Gradient with Momentum (APG-M)}\label{alg:apgm}
\begin{algorithmic}[1]
\STATE \textbf{Input:} stepsize $\eta\in(0,1/L]$, initial point $x_0$, set $x_{-1}=x_0$, $t_0=1$.
\FOR{$k=0,1,2,\dots$}
    \STATE $t_{k+1}\gets \dfrac{1+\sqrt{1+4t_k^{2}}}{2}$          \COMMENT{update momentum scalar}
    \STATE $\beta_k\gets \dfrac{t_k-1}{t_{k+1}}$                    \COMMENT{momentum weight}
    \STATE $y_k \gets x_k + \beta_k (x_k-x_{k-1})$                \COMMENT{extrapolation}
    \STATE $x_{k+1}\gets \prox_{\eta g}\!\bigl( y_k - \eta \nabla f(y_k) \bigr)$
    \STATE \textbf{optional:} check stopping criterion (e.g., $\norm{x_{k+1}-x_k}\le \epsilon$)
\ENDFOR
\end{algorithmic}
\end{algorithm}

\section*{Dry Run Example}
We minimise the scalar problem  

\[
F(w)=f(w)+g(w),\qquad f(w)=(w-3)^2,\qquad g\equiv0 .
\]

Thus $\nabla f(w)=2(w-3)$ and $\prox_{\eta g}(v)=v$.
Choose $\eta=0.1$, initialise $w_0=0$, and use the FISTA schedule.

\begin{center}
\begin{tabular}{c|c|c|c|c}
$k$ & $t_k$ & $\beta_{k-1}$ & $y_{k-1}$ & $w_k$ (after update)\\ \hline
0 & $1$ & $-$ & $-$ & $w_0=0$ \\
1 & $t_1=\frac{1+\sqrt{5}}{2}\approx1.618$ & $\beta_0=\frac{t_0-1}{t_1}=0$ & $y_0=w_0+0\cdot(w_0-w_{-1})=0$ &
\begin{aligned}
&\nabla f(y_0)=2(0-3)=-6\\
&\text{gradient step: }y_0-\eta\nabla f(y_0)=0-0.1(-6)=0.6\\
& w_1=0.6
\end{aligned}\\ \hline
2 & $t_2=\frac{1+\sqrt{1+4t_1^{2}}}{2}\approx2.193$ & $\beta_1=\frac{t_1-1}{t_2}\approx0.282$ &
$y_1=w_1+\beta_1(w_1-w_0)=0.6+0.282(0.6-0)=0.7692$ &
\begin{aligned}
&\nabla f(y_1)=2(0.7692-3)=-4.4616\\
&y_1-\eta\nabla f(y_1)=0.7692-0.1(-4.4616)=1.2154\\
& w_2=1.2154
\end{aligned}\\ \hline
3 & $t_3\approx2.749$ & $\beta_2=\frac{t_2-1}{t_3}\approx0.434$ &
$y_2=w_2+\beta_2(w_2-w_1)=1.2154+0.434(1.2154-0.6)=1.473$ &
\begin{aligned}
&\nabla f(y_2)=2(1.473-3)=-3.054\\
&y_2-\eta\nabla f(y_2)=1.473-0.1(-3.054)=1.778\\
& w_3=1.778
\end{aligned}
\end{tabular}
\end{center}

Objective values:  

\[
F(w_0)=9,\;F(w_1)=5.76,\;F(w_2)=3.10,\;F(w_3)=1.58,
\]

illustrating the accelerated decrease.

\newpage
\section*{Assumptions}
\begin{assumption}[Smoothness of $f$]\label{as:lip}
$f$ is convex and differentiable with $L$‑Lipschitz gradient, i.e.  

\[
\forall x,y\in\R^{d}:\quad 
\norm{\nabla f(x)-\nabla f(y)}\le L\norm{x-y}.
\tag{A1}
\]
\end{assumption}

\begin{assumption}[Convexity of $g$]\label{as:convex}
$g$ is proper, closed and convex (possibly non‑smooth).  
Its proximal mapping $\prox_{\eta g}$ is single‑valued for every $\eta>0$.
\tag{A2}
\end{assumption}

\begin{assumption}[Stepsize]\label{as:stepsize}
The stepsize satisfies  

\[
0<\eta\le \frac{1}{L}.
\tag{A3}
\]
\end{assumption}

\begin{assumption}[Momentum schedule]\label{as:momentum}
The sequence $\{t_k\}$ is defined by  

\[
t_{0}=1,\qquad 
t_{k+1}= \frac{1+\sqrt{1+4t_{k}^{2}}}{2},
\]

and $\beta_k=\frac{t_{k}-1}{t_{k+1}}$.  
Then $t_k\ge \frac{k+1}{2}$ and $0\le\beta_k<1$ for all $k\ge0$.
\tag{A4}
\end{assumption}

\section*{Main Convergence Theorem}
\begin{theorem}[Accelerated $O(1/k^{2})$ rate]\label{thm:rate}
Let $\{x_k\}$ be generated by \eqref{APG-M} with the parameters satisfying Assumptions \ref{as:lip}–\ref{as:momentum}.  
Denote $x^{\star}\in\arg\min_{x}F(x)$. Then for every $k\ge 1$  

\[
F(x_k)-F(x^{\star})\;\le\;
\frac{2\,\norm{x_{0}-x^{\star}}^{2}}{\eta \, (k+1)^{2}} .
\tag{1}
\]

Consequently $\displaystyle \lim_{k\to\infty}F(x_k)=F(x^{\star})$ and the convergence
rate is $O(1/k^{2})$.
\end{theorem}

\section*{Theoretical Properties}
\begin{enumerate}[label=\arabic*.]
\item \textbf{Stability.} The choice $\eta\le 1/L$ guarantees the descent property
\[
F(x_{k+1})\le F(y_k)-\frac{1}{2\eta}\norm{x_{k+1}-y_k}^{2}.
\tag{P1}
\]
\item \textbf{Hyper‑parameter sensitivity.} The momentum schedule \eqref{A4} is
parameter‑free; any other non‑increasing $\beta_k\in[0,1)$ that satisfies
$\beta_k\le \frac{k}{k+3}$ yields the same $O(1/k^{2})$ bound (see Lemma~\ref{lem:beta}).
\item \textbf{Comparison.} 
\begin{itemize}
  \item Classical proximal gradient (no momentum) attains $O(1/k)$.
  \item APG‑M improves to $O(1/k^{2})$ without additional oracle calls.
\end{itemize}
\end{enumerate}

\section*{Discussion}
The proof hinges on three standard tools:

\begin{itemize}
\item \emph{Descent Lemma} for $L$‑smooth $f$,
\item \emph{Three‑point inequality} for the proximal operator,
\item A cleverly constructed potential function that telescopes thanks to the
specific choice of $\{t_k\}$.
\end{itemize}

The result holds for any convex $g$, thus covering $\ell_{1}$ regularisation,
indicator functions of convex sets, etc. Extensions to stochastic gradients or
strongly convex $f$ follow by minor modifications of the potential.

\newpage
\section*{Preliminaries (Definitions and Lemmas)}
\begin{lemma}[Descent Lemma]\label{lem:descent}
If $\nabla f$ is $L$‑Lipschitz, then for all $x,y\in\R^{d}$  

\[
f(y)\le f(x)+\langle\nabla f(x),y-x\rangle+\frac{L}{2}\norm{y-x}^{2}.
\tag{L1}
\]
\end{lemma}
\begin{proof}
Standard; follows from integrating the Lipschitz condition.
\end{proof}

\begin{lemma}[Three‑point inequality for $\prox$]\label{lem:prox}
For any $v\in\R^{d}$ and any $u\in\R^{d}$, letting $p:=\prox_{\eta g}(v)$ we have  

\[
g(p)+\frac{1}{2\eta}\norm{p-v}^{2}\le g(u)+\frac{1}{2\eta}\norm{u-v}^{2}
-\frac{1}{2\eta}\norm{p-u}^{2}.
\tag{L2}
\]
\end{lemma}
\begin{proof}
Immediate from the optimality condition
$0\in \partial g(p)+\frac{1}{\eta}(p-v)$ and convexity of $g$.
\end{proof}

\begin{lemma}[Properties of the momentum sequence]\label{lem:beta}
The sequence $\{t_k\}$ defined in Assumption \ref{as:momentum} satisfies  

\[
t_{k+1}^{2}-t_{k+1}\;=\;t_{k}^{2},\qquad 
\beta_{k}=1-\frac{t_{k}}{t_{k+1}}.
\tag{L3}
\]
\end{lemma}
\begin{proof}
Algebraic manipulation of the recursion for $t_{k+1}$ yields $t_{k+1}^{2}=t_{k+1}+t_{k}^{2}$,
which after rearranging gives the first identity. The second follows from the
definition of $\beta_k$.
\end{proof}

\newpage
\section*{Main Proof (Step‑by‑Step Derivation)}
We prove Theorem~\ref{thm:rate}. Throughout, let $x^{\star}\in\arg\min F$ and define the
potential  

\[
\Phi_k:= t_k^{2}\bigl(F(x_k)-F(x^{\star})\bigr)
        +\frac{1}{2\eta}\norm{z_k-x^{\star}}^{2},
\qquad 
z_k:=t_k x_k-(t_k-1)x_{k-1}.
\tag{P0}
\]

We shall show $\Phi_{k+1}\le\Phi_k$ for all $k\ge0$; telescoping then yields the
desired bound.

\subsection*{Step 1: Relating $z_{k+1}$ and $z_k$}
From the definition of $y_k$ and the momentum parameters,
\[
y_k = x_k + \beta_k (x_k-x_{k-1})
    = \frac{t_k-1}{t_{k+1}}x_{k-1} + \frac{t_{k+1}-t_k+1}{t_{k+1}}x_k .
\tag{1}
\]
Using Lemma~\ref{lem:beta} we obtain
\[
z_{k+1}=t_{k+1}x_{k+1}-(t_{k+1}-1)x_k
      = t_{k+1}\bigl(x_{k+1}-y_k\bigr)+z_k .
\tag{2}
\]

\subsection*{Step 2: Apply the three‑point inequality}
Let $v_k:=y_k-\eta\nabla f(y_k)$. By the algorithm,
$x_{k+1}= \prox_{\eta g}(v_k)$.  
Applying Lemma~\ref{lem:prox} with $p=x_{k+1}$, $v=v_k$ and $u=x^{\star}$ gives  

\[
g(x_{k+1})+\frac{1}{2\eta}\norm{x_{k+1}-v_k}^{2}
\le g(x^{\star})+\frac{1}{2\eta}\norm{x^{\star}-v_k}^{2}
-\frac{1}{2\eta}\norm{x_{k+1}-x^{\star}}^{2}.
\tag{3}
\]

\subsection*{Step 3: Bound $f(x_{k+1})$ using the descent lemma}
Since $v_k = y_k-\eta\nabla f(y_k)$, Lemma~\ref{lem:descent} with $x=y_k$ and $y=x_{k+1}$ yields  

\[
f(x_{k+1})\le f(y_k)+\langle\nabla f(y_k),x_{k+1}-y_k\rangle
            +\frac{L}{2}\norm{x_{k+1}-y_k}^{2}.
\tag{4}
\]

Because $\eta\le 1/L$, we have $\frac{L}{2}\le\frac{1}{2\eta}$, thus

\[
f(x_{k+1})\le f(y_k)+\langle\nabla f(y_k),x_{k+1}-y_k\rangle
            +\frac{1}{2\eta}\norm{x_{k+1}-y_k}^{2}.
\tag{5}
\]

\subsection*{Step 4: Combine $f$ and $g$}
Adding inequality \eqref{3} and \eqref{5} we obtain

\[
\begin{aligned}
F(x_{k+1}) 
&\le f(y_k)+g(x^{\star})+\langle\nabla f(y_k),x_{k+1}-y_k\rangle
      +\frac{1}{2\eta}\norm{x_{k+1}-y_k}^{2}\\
&\qquad +\frac{1}{2\eta}\norm{x_{k+1}-v_k}^{2}
      -\frac{1}{2\eta}\norm{x_{k+1}-x^{\star}}^{2}
      +\frac{1}{2\eta}\norm{x^{\star}-v_k}^{2}.
\end{aligned}
\tag{6}
\]

Observe that $v_k=y_k-\eta\nabla f(y_k)$, hence  

\[
x_{k+1}-v_k = (x_{k+1}-y_k)+\eta\nabla f(y_k).
\tag{7}
\]

Expanding the squared norm in \eqref{7} and simplifying the inner products
cancels the term $\langle\nabla f(y_k),x_{k+1}-y_k\rangle$:

\[
\frac{1}{2\eta}\norm{x_{k+1}-v_k}^{2}
= \frac{1}{2\eta}\norm{x_{k+1}-y_k}^{2}
  +\langle\nabla f(y_k),x_{k+1}-y_k\rangle
  +\frac{\eta}{2}\norm{\nabla f(y_k)}^{2}.
\tag{8}
\]

Substituting \eqref{8} into \eqref{6} eliminates the linear term and yields

\[
F(x_{k+1})\le f(y_k)+g(x^{\star})
            +\frac{1}{\eta}\norm{x_{k+1}-y_k}^{2}
            +\frac{\eta}{2}\norm{\nabla f(y_k)}^{2}
            -\frac{1}{2\eta}\norm{x_{k+1}-x^{\star}}^{2}
            +\frac{1}{2\eta}\norm{x^{\star}-v_k}^{2}.
\tag{9}
\]

\subsection*{Step 5: Express $\norm{x^{\star}-v_k}^{2}$}
Since $v_k=y_k-\eta\nabla f(y_k)$,

\[
\norm{x^{\star}-v_k}^{2}
= \norm{x^{\star}-y_k}^{2}
  -2\eta\langle\nabla f(y_k),x^{\star}-y_k\rangle
  +\eta^{2}\norm{\nabla f(y_k)}^{2}.
\tag{10}
\]

Insert \eqref{10} into \eqref{9} and regroup terms:

\[
\begin{aligned}
F(x_{k+1})-F(x^{\star})
&\le \underbrace{f(y_k)-f(x^{\star})}_{\text{convexity}} 
   +\frac{1}{\eta}\norm{x_{k+1}-y_k}^{2}
   -\frac{1}{2\eta}\norm{x_{k+1}-x^{\star}}^{2}\\
&\quad +\frac{1}{2\eta}\norm{x^{\star}-y_k}^{2}
   -\langle\nabla f(y_k),x^{\star}-y_k\rangle
   +\frac{\eta}{2}\norm{\nabla f(y_k)}^{2}.
\end{aligned}
\tag{11}
\]

By convexity of $f$, $f(y_k)-f(x^{\star})\le\langle\nabla f(y_k),y_k-x^{\star}\rangle$,
which cancels the linear term in \eqref{11}. Consequently,

\[
F(x_{k+1})-F(x^{\star})
\le \frac{1}{2\eta}\bigl(\norm{x^{\star}-y_k}^{2}
      -\norm{x_{k+1}-x^{\star}}^{2}\bigr)
   +\frac{1}{\eta}\norm{x_{k+1}-y_k}^{2}
   +\frac{\eta}{2}\norm{\nabla f(y_k)}^{2}.
\tag{12}
\]

\subsection*{Step 6: Introduce the potential $\Phi_k$}
Multiply \eqref{12} by $t_{k+1}^{2}$ and add the term
$\frac{1}{2\eta}\norm{z_{k+1}-x^{\star}}^{2}$.
Using the relation \eqref{2} for $z_{k+1}$ we obtain after straightforward algebra

\[
\Phi_{k+1}
\le \Phi_k
   -\frac{t_{k+1}}{2\eta}\norm{x_{k+1}-y_k}^{2}
   -\frac{t_{k+1}^{2}\eta}{2}\norm{\nabla f(y_k)}^{2}.
\tag{13}
\]

Both subtracted terms are non‑negative; therefore

\[
\Phi_{k+1}\le\Phi_k\qquad\forall k\ge0.
\tag{14}
\]

\subsection*{Step 7: Telescoping}
Iterating \eqref{14} gives $\Phi_{k}\le\Phi_{0}$.
Since $z_0 = x_0$ and $t_0=1$,  

\[
\Phi_0 = \frac{1}{2\eta}\norm{x_0-x^{\star}}^{2}.
\tag{15}
\]

Discarding the non‑negative proximal term in $\Phi_k$ yields

\[
t_k^{2}\bigl(F(x_k)-F(x^{\star})\bigr)
\le \frac{1}{2\eta}\norm{x_0-x^{\star}}^{2}.
\tag{16}
\]

Finally, using $t_k\ge\frac{k+1}{2}$ from Lemma~\ref{lem:beta} we arrive at

\[
F(x_k)-F(x^{\star})\le
\frac{2\,\norm{x_0-x^{\star}}^{2}}{\eta (k+1)^{2}},
\]
which is exactly \eqref{1}. ∎

\section*{Rate Derivation (With Constants)}
The constant in \eqref{1} can be sharpened by keeping the exact bound
$t_k\ge \frac{k+1}{2}$ (equality for the FISTA schedule). Hence the optimal
asymptotic constant is $2/\eta$.

\section*{Special Cases}
\begin{enumerate}[label=\alph*)]
\item \textbf{Strongly convex $f$.}  
If $f$ is $\mu$‑strongly convex, the same proof yields a linear rate
$F(x_k)-F(x^{\star})\le C\bigl(1-\sqrt{\mu\eta}\bigr)^{k}$.
\item \textbf{Pure gradient descent ($g\equiv0$).}  
The algorithm reduces to Nesterov’s accelerated gradient method,
and the bound \eqref{1} coincides with the classical $O(1/k^{2})$ result.
\item \textbf{Indicator $g=\iota_{\mathcal{C}}$.}  
When $g$ is the indicator of a convex set $\mathcal{C}$, $\prox_{\eta g}$ is the Euclidean
projection onto $\mathcal{C}$, and the theorem guarantees accelerated convergence of
projected gradient descent.
\end{enumerate}

\end{document}